\chapter*{Словарь терминов}             % Заголовок
\addcontentsline{toc}{chapter}{Словарь терминов}  % Добавляем его в оглавление

{\bfseries\boldmath ab initio} : квантово-химический расчет из первых принципов, выполняемый без подгонки параметров непосредственно под анализируемый экспериментальный спектр.

{\bfseries\boldmath Bruker IFS~125HR} : Фурье-спектрометр высокого разрешения, использованный для регистрации экспериментальных инфракрасных спектров.

{\bfseries\boldmath \(C_{2v}\)} : точечная группа симметрии молекулы сероводорода H$_2$S и симметричных изотопологов типа D$_2$S.

{\bfseries\boldmath \(C_s\)} : точечная группа симметрии, содержащая одну плоскость симметрии; в работе используется при описании молекулы HDS.

{\bfseries\boldmath \(d_{\mathrm{rms}}\)} : среднеквадратичное отклонение экспериментальных значений от рассчитанных.

{\bfseries\boldmath D$_2$S} : дидейтерированный сероводород, изотополог H$_2$S, в котором оба атома водорода заменены атомами дейтерия.

{\bfseries\boldmath GF} : метод Вильсона для расчета нормальных колебаний молекул; $G$ обозначает матрицу кинематических коэффициентов, $F$ --- силовую матрицу.

{\bfseries\boldmath GeH$_4$} : герман, тетрагидрид германия.

{\bfseries\boldmath HDS} : монодейтерированный сероводород; в работе также рассматриваются изотопические разновидности HD$^{32}$S и HD$^{34}$S.

{\bfseries\boldmath ИК} : инфракрасная область спектра; в тексте используется в выражениях <<ИК-спектры>> и <<ИК-диапазон>>.

{\bfseries\boldmath KBr} : бромид калия, материал оптических элементов, указанный в условиях регистрации спектров.

{\bfseries\boldmath MCT} : детектор на основе теллурида ртути--кадмия (mercury cadmium telluride), применяемый при регистрации инфракрасных спектров.

{\bfseries\boldmath SiH$_4$} : силан, тетрагидрид кремния.

{\bfseries\boldmath STDS} : Spherical Top Data System, база данных и программный комплекс для расчета и анализа спектров молекул типа сферического волчка.

{\bfseries\boldmath \(T_d\)} : точечная группа тетраэдрической симметрии, используемая при описании молекул типа XY$_4$.

{\bfseries\boldmath XY$_2$} : обобщенная формула трехатомной молекулы с центральным атомом X и двумя атомами Y.

{\bfseries\boldmath XY$_4$} : обобщенная формула тетраэдрической молекулы с центральным атомом X и четырьмя атомами Y.
